\section{Traps, exceptions and interrupts}
\textit{Sources: \texttt{design/core/interrupt/}, \texttt{trap\_sequencer/},
\texttt{wb\_trap\_unit/}, \texttt{csr\_unit/}, and their wiring in
\texttt{core\_top.sv}.}

A \emph{trap} is any forced, non-sequential change of control flow into a
handler: a synchronous \emph{exception} caused by the instruction itself
(illegal opcode, \sig{ecall}, a page/access fault, a misaligned access), or an
asynchronous \emph{interrupt} from outside the pipeline (timer, software, or
external). This section explains how ntiny detects, prioritises, commits, and
returns from traps.

\diagram{trap-flow}{Trap flow. Exceptions are detected in the stage that can
see them (top); interrupts arrive from CLINT/PLIC/Sstc. \texttt{interrupt\_ctrl}
picks the winner, its cause, the saved PC (epc) and whether it delegates to
S-mode; \texttt{wb\_trap\_unit} commits it atomically at writeback, which drives
the CSR updates, the pipeline flush, and the redirect to the handler.}{fig:trap-flow}{0.78}

\subsection{The big idea: traps commit at writeback (IWB)}
ntiny resolves every trap, interrupt and return (\sig{mret}/\sig{sret}/\sig{dret})
\textbf{atomically at the IWB stage} (\sig{wb\_trap\_unit}). An instruction is
tagged during decode/execute with a \sig{wb\_event} (\sig{WB\_TRAP},
\sig{WB\_XRET}, \sig{WB\_DRET}); when that tagged ``carrier'' reaches IWB, the
trap fires. Committing at one fixed point---rather than wherever the fault was
first noticed---is what makes the privileged state consistent: it removes the
old races between an \sig{mret} and the \sig{csrrw mepc} just ahead of it, and
between a page-table walk and the load that triggered it. When a trap fires,
\sig{wb\_trap\_unit} asserts \sig{kill\_iwb}: the carrier does \emph{not} write
its register or commit its store, and a four-stage flush (IE/IMEM/IWB plus the
fetch redirect) kills every younger wrong-path instruction.

\subsection{Exception sources}
Each exception is detected in the earliest stage that has the information, then
carried to the commit point. The saved PC (\sig{epc}) is the faulting
instruction for exceptions (so the handler can fix up and retry), and
\sig{mtval}/\sig{stval} carries the offending address or instruction word.

\begin{center}\small
\begin{tabularx}{\textwidth}{l c l Y}
\toprule
\textbf{exception} & \textbf{cause} & \textbf{detected in} & \textbf{tval} \\
\midrule
Instruction access fault (PMP) & 1  & IF (PMP / PTW) & faulting VA (or straddled half) \\
Illegal instruction            & 2  & ID, or IE (bad CSR) & instruction word \\
Breakpoint (\sig{ebreak})       & 3  & ID & 0 \\
Load address misaligned        & 4  & IE & address \\
Load access fault (PMP)        & 5  & IE & address \\
Store/AMO misaligned           & 6  & IE & address \\
Store/AMO access fault (PMP)   & 7  & IE & address \\
ECALL from U / S / M           & 8/9/11 & ID & 0 \\
Instruction page fault         & 12 & IF (ITLB / PTW) & faulting VA \\
Load page fault                & 13 & IE (DTLB / PTW) & address \\
Store/AMO page fault           & 15 & IE (DTLB / PTW) & address \\
\bottomrule
\end{tabularx}
\end{center}

Two implementation details worth knowing. First, instruction-side faults
(causes 1 and 12) take \emph{two} paths: a combinational one that rides through
the \sig{fetch\_buffer} on the same entry as the faulting word (so it fires when
that word reaches the buffer head), and a direct ``PTW-completion'' path for
faults discovered by the page-table walker (which has no live fetch request to
ride). Second, the data-side fault's \emph{is-store} flag is registered
(\sig{trap\_sequencer}) because the trap can fire a cycle or two after the IE
stage has already flushed and dropped the live store signal.

\subsection{Interrupt sources}
Three interrupt lines arrive from the uncore and become pending bits in
\sig{mip}: machine timer (\sig{MTIP}, from the CLINT \sig{mtime}\,$\ge$\,\sig{mtimecmp}
compare), machine software (\sig{MSIP}, CLINT), and machine external (\sig{MEIP},
PLIC). Their S-mode counterparts (\sig{STIP}/\sig{SSIP}/\sig{SEIP}) are either
injected by M-mode software via \sig{mip}, sourced from the PLIC's S-mode
context, or---for the timer---generated in hardware by the \textbf{Sstc}
extension (\sig{stimecmp} compared against \sig{mtime} when \sig{menvcfg.STCE=1}).
An interrupt is \emph{taken} only when its pending bit and its enable bit
(\sig{mie}/\sig{sie}) are both set \emph{and} the global enable for the current
privilege allows it: in M-mode that is \sig{mstatus.MIE}; below M-mode, M
interrupts are always enabled. Delegation (\sig{mideleg}) sends a delegated
interrupt to S-mode instead of M-mode when not already in M-mode.

\latencynote{Interrupts are held off until the pipeline actually has an
in-flight instruction to pin them to (\sig{pc\_id} or \sig{pc\_ie} non-zero). If
a timer fires into an empty/draining pipeline, naively using the fetch-stage PC
as \sig{epc} would set the return address \emph{ahead} of instructions still in
flight, double-executing or skipping them on \sig{sret}. Deferring the interrupt
one cycle until the pipeline is populated avoids that.}

\subsection{Priority}
Synchronous exceptions outrank interrupts on the same instruction, and older
(deeper-in-the-pipeline) exceptions outrank younger ones. The order, highest
first: IE-stage data faults (PMP access, then page fault, then misalign, then
illegal-CSR) $\rightarrow$ IF/ID-stage faults (instruction access, instruction
page, illegal, \sig{ecall}, \sig{ebreak}) $\rightarrow$ interrupts (external,
then software, then timer; M-mode before S-mode). At the IWB commit, a debug
entry outranks everything, then an interrupt preempts the carrier, then a
synchronous trap, then \sig{xret}.

\subsection{Choosing the saved PC for an interrupt}
The trickiest part of the trap logic---and historically the most bug-prone---is
picking \sig{epc} for an \emph{asynchronous} interrupt, because the ``current''
instruction depends on what the pipeline is doing that cycle. \sig{interrupt\_ctrl}
selects \sig{pc\_for\_async} with a priority fallback:
\begin{enumerate}[leftmargin=1.4em,itemsep=1pt]
\item a branch/jump is committing this cycle $\rightarrow$ its architectural
  next-PC (\sig{branch\_recovery\_target}), which collapses all
  predicted$\times$actual direction cases into the correct resume point;
\item an IE-stage op has not committed (AMO in progress, data MMU stall)
  $\rightarrow$ \sig{pc\_ie} (re-execute it after \sig{sret});
\item \sig{pc\_id} is non-zero $\rightarrow$ \sig{pc\_id} (the next instruction
  to enter IE);
\item \sig{pc\_ie} is non-zero $\rightarrow$ \sig{pc\_ie + insn\_len}
  (length-aware: 2 for compressed, 4 otherwise);
\item otherwise the fetch-stage PC (only when the pipeline is empty, e.g.
  just after an \sig{xret}).
\end{enumerate}
Several past Linux-boot hangs were exactly a wrong choice here (e.g.\ a timer
firing the same cycle as a taken-and-correctly-predicted branch, where the naive
fall-through PC is the wrong resume target).

\subsection{Trap entry: what hardware writes}
When the trap commits, \sig{csr\_unit} performs, in one cycle, all of:
\begin{itemize}[leftmargin=1.4em,itemsep=1pt]
\item \sig{mepc}/\sig{sepc} $\leftarrow$ the selected \sig{epc};
\item \sig{mcause}/\sig{scause} $\leftarrow$ the cause (interrupt bit set for
  asynchronous);
\item \sig{mtval}/\sig{stval} $\leftarrow$ the trap value (0 for interrupts and
  \sig{ebreak});
\item \sig{mstatus}: for an M-trap, \sig{MPP}$\leftarrow$old privilege,
  \sig{MPIE}$\leftarrow$\sig{MIE}, \sig{MIE}$\leftarrow$0 (the S-trap form uses
  \sig{SPP}/\sig{SPIE}/\sig{SIE}, with \sig{SPP} a single bit);
\item \sig{priv\_level} $\leftarrow$ M (\sig{2'b11}) or S (\sig{2'b01})
  per delegation.
\end{itemize}
The handler address is \sig{mtvec}/\sig{stvec} with its low two bits stripped;
in vectored mode (base[0]=1) an \emph{interrupt} targets base\,$+\,4\times$cause,
while exceptions always go to the base.

\subsection{The carrier still commits its CSR write}\label{sec:traps-carrier}
\latencynote{Subtle but important: when an asynchronous interrupt preempts a CSR
instruction (say \sig{csrsi sstatus,2} that sets \sig{SIE}), that CSR write must
still take effect, because the chosen \sig{epc} points \emph{past} it---the
architecture has logically retired it. \sig{csr\_unit} therefore computes the
new \sig{mstatus} by first applying the in-flight CSR write and \emph{then}
layering the trap-entry overrides (\sig{MPIE}$\leftarrow$post-write \sig{MIE},
etc.) on top. A real Linux scheduler hang was caused by the earlier code
suppressing the carrier's CSR write on a coincident trap: \sig{SIE} stayed 0, so
the kernel resumed with interrupts disabled and warned in \texttt{do\_idle()}.}

\subsection{Trap return: \texttt{mret} / \texttt{sret} / \texttt{dret}}
A return is decoded, tagged \sig{WB\_XRET} (or \sig{WB\_DRET}), and commits at
IWB. \sig{mret} restores \sig{MIE}$\leftarrow$\sig{MPIE}, sets \sig{MPIE}=1,
\sig{MPP}=U, clears \sig{MPRV} if returning below M, and jumps to \sig{mepc};
\sig{priv\_level}$\leftarrow$\sig{MPP}. \sig{sret} is the analogous S-mode form
(\sig{SIE}$\leftarrow$\sig{SPIE}, \sig{SPIE}=1, \sig{SPP}=U, jump to \sig{sepc}).
\sig{dret} restores \sig{dpc} and leaves debug mode without a privilege change.

Because a return changes privilege, the fetch of the first handler-exit
instruction must use the \emph{new} privilege for translation. A small fetch FSM
(\sig{xret\_fetch\_ctrl}) drains speculative wrong-privilege fetches, waits one
cycle for \sig{priv\_level} to settle, then re-fetches from the return target. A
companion guard in \sig{trap\_sequencer} suppresses stale instruction faults from
the \emph{old} privilege's in-flight fetch while the \sig{sret} CSR side effects
land---otherwise a cold-ITLB page walk on the \sig{sret} target could, several
cycles later, overwrite \sig{sepc}/\sig{scause} with a garbage fault.

\subsection{Interaction with the stall chain}
\latencynote{Traps share the pipeline with the variable-latency hazards from the
rest of this document. Two cases matter. (1) An interrupt coincident with a
\emph{stalled load} must save \sig{pc\_ie} (case 2 above) so the load
re-executes after \sig{sret}; if its data happens to arrive the same cycle, the
IE flush suppresses the (now wrong-path) register write, and the post-\sig{sret}
re-execution re-reads memory correctly. (2) A speculative fetch past an
unconditional jump can sit in a no-execute PMP region; its access fault must be
\emph{squashed} when an older redirect (branch/\sig{xret}/debug) is firing, so a
wrong-path fault never reaches the handler. Both are handled by gating the fault
on the live redirect/flush signals rather than letting it fire blindly.}
